\subsection{欧拉路径 (Eulerian Path)}
\paragraph{有向图}
\begin{lstlisting}
vi e[100010];
ll cur[100010];
ll in[100010],out[100010];
ll ans[1000010],c;
void dfs(ll u){
	for(int i=cur[u];i<e[u].size();i=cur[u]){
		ll v=e[u][i];cur[u]=i+1;
		dfs(v);
	}ans[++c]=u;
}
int main(){
	n=d,m=d;c=0;
    f(i,1,n)e[i].clear(),cur[i]=in[i]=out[i]=0;
    f(i,1,m){
		ll u=d,v=d;
		e[u].pb(v);
		in[v]++,out[u]++;
	}f(i,1,n)sort(e[i].begin(),e[i].end());
	ll x = -1, y = -1; // x: start, y: end
	f(i,1,n){
		// if(i==1)cout<<in[i]<<" "<<out[i]<<endl;
		if(in[i]==out[i])continue;
		if(abs(in[i]-out[i])>1){pn;return 0;}
		if(in[i]==out[i]+1){if(y!=-1){pn;return 0;}y=i;}
		if(in[i]==out[i]-1){if(x!=-1){pn;return 0;}x=i;}
	}if(x==-1){x=1;f(i,1,n)if(out[i]){x=i;break;}}dfs(x);
    if(c!=m+1){pn;return 0;}
	for(int i=c;i>=1;i--)cout<<ans[i]<<" ";
	return 0;
}
\end{lstlisting}

\paragraph{无向图}
无向图：把经过过的边标记一下不要走即可 (Hierholzer 算法)。

\paragraph{多测与注意} 保留递归 Hierholzer，已补非零度起点与用尽所有边的检查。多测初始化已放在读边前；边很多时 DFS 栈深为边数。
